\begin{table}[t]
\caption{
In additional to the standard Linux Bash commands, we provide SWE-agent with specialized tools, including an interactive file viewer, search functionalities, and edit tools for the open file. Required arguments are enclosed in \texttt{<>} and optional arguments are in \texttt{[]}. The last column shows the documentation presented to the LM.
}
\label{tab:commands}
\vspace{0.25em}
\centering
\begin{tabular}{p{0.08\linewidth}>{\raggedright\arraybackslash}p{0.33\linewidth}p{0.48\linewidth}}
\toprule
\multicolumn{1}{l}{Category}& \multicolumn{1}{c}{Command} & \multicolumn{1}{c}{Documentation} \\
\midrule
% \multicolumn{2}{l}{\textit{File viewer}} \\
% \cmidrule(lr){1-2}
\textit{File viewer} & \textbf{open}\texttt{ <path> [<line\_number>]} & Opens the file at the given path in the editor. If \texttt{line\_number} is provided, the window will move to include that line. \\
\addlinespace
& \textbf{goto}\texttt{ <line\_number>} & Moves the window to show line\_number. \\
\addlinespace
& \textbf{scroll\_down} & Moves the window up 100 lines. \\
\addlinespace
& \textbf{scroll\_up} & Moves the window down 100 lines. \\
% \cmidrule(lr){1-2}
% \multicolumn{2}{l}{\textit{Search tools}} \\
% \cmidrule(lr){1-2}
\cmidrule(lr){1-3}
\textit{Search tools} & \textbf{search\_file}\texttt{ <search\_term> [<file>]} & Searches for \texttt{search\_term} in file. If file is not provided, searches in the current open file. \\
\addlinespace
& \textbf{search\_dir}\texttt{ <search\_term> [<dir>]} & Searches for \texttt{search\_term} in all files in dir. If dir is not provided, searches in the current directory. \\
\addlinespace
& \textbf{find\_file}\texttt{ <file\_name> [<dir>]} & Finds all files with the given name in dir. If dir is not provided, searches in the current directory. \\
% \cmidrule(lr){1-2}
% \multicolumn{2}{l}{\textit{Editing}} \\
% \cmidrule(lr){1-2}
\cmidrule(lr){1-3}
\textit{File \mbox{editing}} & {\textbf{edit} \texttt{<n>:<m>}  \texttt{<replacement\_text>} \textbf{end\_of\_edit}} & Replaces lines n through m (inclusive) with the given text in the open file. All of the \texttt{replacement\_text} will be entered, so make sure your indentation is formatted properly. Python files will be checked for syntax errors after the edit. If an error is found, the edit will not be executed. Reading the error message and modifying your command is recommended as issuing the same command will return the same error. \\
\addlinespace
& \textbf{create}\texttt{ <filename>} & Creates and opens a new file with the given name. \\
% \cmidrule(lr){1-2}
% \multicolumn{2}{l}{\textit{High-level control}} \\
% \cmidrule(lr){1-2}
\cmidrule(lr){1-3}
\textit{Task} & \textbf{submit} & Generates and submits the patch from all previous edits and closes the shell. \\
\bottomrule
\end{tabular}
\vspace{-2em}
\end{table}